% ============================================================
% Ethics and Reproducibility statements.  ICLR excludes both from the 9-page
% main-text limit, so nothing here competes with the body for space.  They are
% placed after the conclusion and before the references, per the ICLR template.
% ============================================================

\section*{Ethics Statement}

\textbf{Data.}  Every experiment uses public benchmark data and involves no human
subjects and no personally identifying information.  The forecasting datasets
(ETTh1, ETTh2, ETTm2, Weather and Electricity) are standard public
releases of instrument and utility measurements.  The ILI cell uses public
aggregate influenza-surveillance counts, not individual health records.  The IoT
negative control uses N-BaIoT, network traffic captured from a laboratory
testbed of consumer devices rather than from a deployed network.  We introduce no
new data collection.

\textbf{The main risk in this paper is that its negative result is read too
broadly.}  We report that on these benchmarks we could not construct a cell in
which fine-tuning demonstrably damages a pre-trained capability, and that a
diagnostic we proposed for predicting such damage has no positive instance to
predict.  That is a statement about six standard forecasting benchmarks and three
backbones, and \S\ref{sec:limitations} states what it does not license.  It is
\emph{not} evidence that fine-tuning is safe in general, and a practitioner who
reads it as permission to skip the frozen-encoder control would be drawing the
opposite of our recommendation: our prescription is to \emph{run} that control,
because we found the observational shortcuts unreliable.  Domains where a
pre-trained checkpoint holds a large, well-established advantage --- which these
benchmarks do not, on 27 of the 32 cells we screened --- are exactly where the
question remains open.

\textbf{Research integrity.}  This version corrects an earlier version of the
same work in two ways that changed its conclusions rather than its presentation,
and we flag both rather than absorbing them silently.  An earlier version's value
gate used an unfitted per-window trend extrapolation as its linear baseline; a
constant predictor beats that baseline, and correcting it changes the gate status
of 17 of the 21 Moirai cells and removes every degradation instance the paper had
--- that is, the entire class of cell the diagnostic was built to
predict.  Separately, we retract a linear-probe diagnostic
reported earlier: the statistic it rests on is non-zero only where the probe
fails to generalise, and Appendix~\ref{app:orthogonal} reports the measurement
that establishes this on the same encoders.  Neither correction required
re-running a model; both are documented with their audit trails in
Appendix~\ref{sec:retraction} and Appendix~\ref{app:gatecorrection}.  Where a
number in an earlier version could not be traced to a run record, we say so at
the site rather than reprinting it.

\textbf{Compute.}  The experiments are small by contemporary standards: no run
exceeds 10{,}000 training windows, and the sweep reported here runs on a single
GPU or on Apple silicon.  We report the protocol deviations forced by the
hardware available (Appendix~\ref{app:layerunfreeze}) rather than presenting the
runs as environment-independent.

\section*{Reproducibility Statement}

The paper is organised so that every quantity can be traced to a stored run
record, and we have tried to make the failure modes we hit ourselves cheap for
others to avoid.

\textbf{Where the numbers come from.}  The cell-level results are not assembled by
hand: \texttt{scripts/cell\_matrix.py} reads the per-run JSON records and emits
the main tables, so a discrepancy between the text and the runs is detectable by
regenerating them.  Each appendix table names the result directory it is read
from, and the cross-cell correlations of Appendix~\ref{app:crosscell} are the
output of named functions in that script rather than transcribed values.  The
per-run records carry the fields the tables use --- \texttt{final\_cka},
\texttt{forgetting\_pct}, \texttt{final\_val\_mse}, \texttt{final\_weight\_drift}
and the early-stopping block --- together with the seed, epoch budget and sample
count, so the dispersion conventions can be recomputed rather than taken on
trust.

\textbf{Two things are required to reproduce the Moirai runs at all.}  First, the
\texttt{uni2ts} \texttt{PackedStdScaler} performs an in-place write that silently
detaches the gradient tape for losses depending on scaled statistics; without the
patch of Appendix~\ref{app:scaler_patch} training can appear to run while not
converging.  Second, the frozen-encoder condition is not a strict freeze --- the
input projections and mask encoding remain trainable, which is why its CKA is not
$1.000$ --- and the strictly-frozen variant is reported separately
(Appendix~\ref{app:strictfreeze}).  Comparisons against a ``frozen'' baseline are
not comparable across papers unless which parameters were frozen is stated.

\textbf{Scoring protocol.}  All headline comparisons are scored on windows
disjoint from those used for early stopping.  This is not a refinement:
\S\ref{sec:analysis:heldout} shows that scoring on the selection windows instead
reverses the sign of the frozen-encoder comparison in 4 of the 31 intervention
cells, always in the same direction.  Reproducing our numbers therefore requires
reproducing the split discipline, and Appendix
Table~\ref{tab:heldout_decomp} gives the per-arm decomposition needed to check
which split a given number was scored on.

\textbf{Pre-registration.}  The prospective test of
\S\ref{sec:analysis:prospective} was registered before any outcome run for
those cells existed.  The predictions were written to
\texttt{preregistration.json} under \texttt{results/}, and committed before the
first condition-B or condition-D run of that arm was launched; the scorer,
\texttt{score\_prospective.py}, reads that file and the finished runs and makes
no decisions of its own.  Both are in the release, and the commit that added
the registration precedes the runs it scores, so the ordering is checkable from
the history rather than asserted by us.

\textbf{What we release.}  The analysis and experiment scripts, the per-run JSON
records behind every table, the pre-registration file, and the \LaTeX{} source
and style files needed to build this paper.  Model checkpoints are excluded for
size; the appendices state which runs retained a checkpoint, because two
re-measurements we would otherwise have made were blocked by runs whose
checkpoints were not kept (Appendix~\ref{app:noisefloor}).
